\documentclass[11pt,a4paper]{article}

\usepackage[utf8]{inputenc}
\usepackage[T1]{fontenc}
\usepackage[spanish,es-nodecimaldot]{babel}
\usepackage{lmodern}
\usepackage{amsmath,amssymb,mathtools}
\usepackage{booktabs}
\usepackage{geometry}
\usepackage{xcolor}
\usepackage{hyperref}
\usepackage{enumitem}
\usepackage{microtype}

\geometry{margin=2.5cm}
\hypersetup{
  colorlinks=true,
  linkcolor=blue!55!black,
  urlcolor=blue!55!black,
  citecolor=blue!55!black
}

\setlength{\parindent}{0pt}
\setlength{\parskip}{0.7em}
\setlist[itemize]{topsep=0.2em,itemsep=0.2em}
\setlist[enumerate]{topsep=0.2em,itemsep=0.2em}

\definecolor{pucpblue}{RGB}{27,65,119}
\definecolor{softblue}{RGB}{237,243,252}
\definecolor{softred}{RGB}{253,238,241}
\definecolor{softgreen}{RGB}{238,247,241}

\newcommand{\Zauto}{\mathcal{Z}^{\mathrm{auto}}_c}
\newcommand{\one}{\mathbf{1}}
\newcommand{\Eauto}{\mathbb{E}_{\Zauto}}
\newcommand{\Covauto}{\operatorname{Cov}_{\Zauto}}
\newcommand{\Devs}{\operatorname{Devs}}
\newcommand{\Uprog}{U_{\mathrm{prog}}}
\newcommand{\Uout}{U_{\mathrm{out}}}
\newcommand{\dd}{\,dz}
\newcommand{\nota}[1]{%
  \begin{center}
  \fcolorbox{pucpblue}{softblue}{%
  \begin{minipage}{0.93\linewidth}
  \textbf{Intuición.} #1
  \end{minipage}}
  \end{center}
}
\newcommand{\supuesto}[1]{%
  \begin{center}
  \fcolorbox{red!60!black}{softred}{%
  \begin{minipage}{0.93\linewidth}
  \textbf{Supuesto no trivial.} #1
  \end{minipage}}
  \end{center}
}
\newcommand{\resultado}[1]{%
  \begin{center}
  \fcolorbox{green!45!black}{softgreen}{%
  \begin{minipage}{0.93\linewidth}
  \textbf{Resultado.} #1
  \end{minipage}}
  \end{center}
}

\title{\textbf{Explicación detallada del modelo del Apéndice B}\\
\large Impacto de la disponibilidad de ChatGPT en desarrolladores activos de GitHub}
\author{Documento de apoyo matemático}
\date{Lima, 2026}

\begin{document}

\maketitle
\tableofcontents
\newpage

\section{Qué intenta hacer el modelo}

El modelo quiere conectar una idea tecnológica con una ecuación estimable. La idea tecnológica es:

\begin{quote}
ChatGPT no mejora todas las tareas por igual. Mejora más aquellas tareas de programación donde su productividad supera a la productividad humana. Si esas mejoras aumentan la productividad esperada de programar, algunas personas que antes no entraban al mercado ahora sí entran.
\end{quote}

La cadena completa es:
\[
\begin{gathered}
\text{mejora por tarea}
\Longrightarrow
\text{ganancia individual}
\Longrightarrow
\text{menor umbral de entrada}\\
\Longrightarrow
\text{más desarrolladores activos}
\Longrightarrow
\text{ecuación de panel}.
\end{gathered}
\]

Este documento desarrolla esa cadena sin saltar pasos. La derivación tiene dos tipos de operaciones:

\begin{itemize}
  \item \textbf{Identidades exactas}: son reescrituras algebraicas. No agregan supuestos.
  \item \textbf{Aproximaciones o supuestos}: simplifican la expresión para obtener una forma interpretable y estimable.
\end{itemize}

Una forma sencilla de recordar el resultado central es:
\[
\Delta Y_{ic} \approx \phi_i \theta_c.
\]

Esto dice que el aumento de desarrolladores por cada 100,000 habitantes en el país \(i\) y lenguaje \(c\) depende de:

\begin{itemize}
  \item \(\theta_c\): exposición del lenguaje a ChatGPT.
  \item \(\phi_i\): capacidad de absorción del país.
\end{itemize}

\section{Notación básica}

\begin{center}
\begin{tabular}{ll}
\toprule
Símbolo & Significado \\
\midrule
\(z \in [0,1]\) & Tarea elemental de programación. \\
\(c\) & Lenguaje de programación. \\
\(i\) & País. \\
\(j\) & Desarrollador o persona potencialmente desarrolladora. \\
\(h_c(z)>0\) & Productividad humana en la tarea \(z\) del lenguaje \(c\). \\
\(a_c(z)\geq 0\) & Productividad de ChatGPT en la tarea \(z\). \\
\(\mu_j \in \{0,1\}\) & Acceso individual a ChatGPT. \\
\(y_{jc}(\mu_j)\) & Producción total del individuo \(j\) en lenguaje \(c\). \\
\(\Zauto\) & Conjunto de tareas donde ChatGPT supera al humano. \\
\(\tau_c\) & Fracción de tareas automatizables. \\
\(s_c(z)\) & Ahorro relativo de ChatGPT en la tarea \(z\). \\
\(\rho_c\) & Ahorro relativo promedio en tareas automatizables. \\
\(\eta\) & Habilidad individual para programar. \\
\(r_i\) & Utilidad de reserva o alternativa fuera de programar. \\
\(\phi_i\) & Capacidad de absorción tecnológica del país \(i\). \\
\(\theta_c\) & Exposición del lenguaje \(c\) a ChatGPT. \\
\(Y_{ict}\) & Desarrolladores activos por 100,000 habitantes. \\
\bottomrule
\end{tabular}
\end{center}

\section{Paso 1: función de producción a nivel de tarea}

La producción en una tarea se define como:
\[
y_{jc}(z;\mu_j)=\max\{h_c(z),\mu_j a_c(z)\}.
\]

Leamos cada parte:

\begin{itemize}
  \item \(h_c(z)\) es lo que produce el humano en la tarea \(z\).
  \item \(a_c(z)\) es lo que produce ChatGPT en la tarea \(z\).
  \item \(\mu_j\) indica si el individuo tiene acceso a ChatGPT.
  \item El operador \(\max\{\cdot,\cdot\}\) significa que se escoge el método más productivo.
\end{itemize}

Si \(\mu_j=0\), no hay acceso a ChatGPT. Entonces:
\[
\mu_j a_c(z)=0\cdot a_c(z)=0.
\]

La producción por tarea queda:
\[
y_{jc}(z;0)=\max\{h_c(z),0\}.
\]

Como se asume \(h_c(z)>0\), el máximo entre \(h_c(z)\) y \(0\) es \(h_c(z)\). Por tanto:
\[
y_{jc}(z;0)=h_c(z).
\]

Si \(\mu_j=1\), sí hay acceso a ChatGPT. Entonces:
\[
\mu_j a_c(z)=1\cdot a_c(z)=a_c(z).
\]

La producción por tarea queda:
\[
y_{jc}(z;1)=\max\{h_c(z),a_c(z)\}.
\]

\subsection{De producción por tarea a producción total}

El intervalo \([0,1]\) representa el continuo de tareas. Integrar sobre \([0,1]\) significa sumar la contribución de todas las tareas:
\[
y_{jc}(\mu_j)=\int_0^1 \max\{h_c(z),\mu_j a_c(z)\}\dd.
\]

Sin IA:
\[
y_{jc}(0)=\int_0^1 h_c(z)\dd.
\]

Con IA:
\[
y_{jc}(1)=\int_0^1 \max\{h_c(z),a_c(z)\}\dd.
\]

\nota{Una integral aquí funciona como una suma continua. Si hubiera 100 tareas discretas, sumaríamos la producción de cada tarea. Como el modelo usa un continuo de tareas entre 0 y 1, usamos una integral.}

\section{Paso 2: conjunto de tareas automatizables}

Definimos:
\[
\Zauto=\{z\in[0,1]:a_c(z)>h_c(z)\}.
\]

Esto significa:
\[
z\in\Zauto
\quad \Longleftrightarrow \quad
a_c(z)>h_c(z).
\]

En palabras:

\begin{quote}
\(\Zauto\) es el conjunto de tareas donde ChatGPT es estrictamente más productivo que el humano.
\end{quote}

El complemento es:
\[
[0,1]\setminus \Zauto.
\]

Esto significa:
\[
z\in[0,1]\setminus \Zauto
\quad \Longleftrightarrow \quad
z\in[0,1]\ \text{y}\ z\notin\Zauto.
\]

Si \(z\notin\Zauto\), entonces no se cumple \(a_c(z)>h_c(z)\). Por tanto:
\[
a_c(z)\leq h_c(z).
\]

Así, el intervalo de tareas se parte en dos grupos:
\[
[0,1]=([0,1]\setminus\Zauto)\cup \Zauto,
\]
con intersección vacía:
\[
([0,1]\setminus\Zauto)\cap\Zauto=\varnothing.
\]

\subsection{Indicadoras}

La función indicadora se define así:
\[
\one[z\in A]=
\begin{cases}
1, & \text{si }z\in A,\\
0, & \text{si }z\notin A.
\end{cases}
\]

Por eso:
\[
f(z)\one[z\in A]=
\begin{cases}
f(z), & \text{si }z\in A,\\
0, & \text{si }z\notin A.
\end{cases}
\]

Al integrar:
\[
\int_0^1 f(z)\one[z\in A]\dd=\int_A f(z)\dd.
\]

La izquierda integra sobre todo \([0,1]\), pero la indicadora apaga la función fuera de \(A\). Por eso es igual a integrar directamente sobre \(A\).

\subsection{Reescritura del máximo}

Para cada tarea \(z\), el máximo puede escribirse como:
\[
\max\{h_c(z),a_c(z)\}
=h_c(z)\one[z\in[0,1]\setminus\Zauto]+a_c(z)\one[z\in\Zauto].
\]

Verificación por casos:

\textbf{Caso 1: \(z\in\Zauto\).} Entonces \(a_c(z)>h_c(z)\). El máximo es \(a_c(z)\). Además:
\[
\one[z\in\Zauto]=1,\qquad
\one[z\in[0,1]\setminus\Zauto]=0.
\]
Por tanto:
\[
h_c(z)\cdot 0+a_c(z)\cdot 1=a_c(z).
\]

\textbf{Caso 2: \(z\in[0,1]\setminus\Zauto\).} Entonces \(a_c(z)\leq h_c(z)\). El máximo es \(h_c(z)\). Además:
\[
\one[z\in\Zauto]=0,\qquad
\one[z\in[0,1]\setminus\Zauto]=1.
\]
Por tanto:
\[
h_c(z)\cdot 1+a_c(z)\cdot 0=h_c(z).
\]

\subsection{Fracción de tareas automatizables}

La fracción de tareas automatizables es:
\[
\tau_c=\int_0^1 \one[a_c(z)>h_c(z)]\dd.
\]

Como \(\one[a_c(z)>h_c(z)]\) vale 1 dentro de \(\Zauto\) y 0 fuera, esto equivale a:
\[
\tau_c=\int_{\Zauto}1\dd.
\]

Entonces \(\tau_c\) mide el tamaño de \(\Zauto\). Como \(\Zauto\subseteq[0,1]\), se cumple:
\[
0\leq\tau_c\leq 1.
\]

\nota{Si \(\tau_c=0.30\), ChatGPT supera al humano en el 30\% de las tareas del lenguaje \(c\).}

\section{Paso 3: ahorro relativo por tarea}

Para \(z\in\Zauto\), se cumple:
\[
a_c(z)>h_c(z)>0.
\]

El ahorro relativo de ChatGPT se define como:
\[
s_c(z)=\frac{a_c(z)-h_c(z)}{h_c(z)}.
\]

Esta expresión mide la ganancia proporcional de ChatGPT respecto al humano. Por ejemplo, si:
\[
h_c(z)=10,\qquad a_c(z)=15,
\]
entonces:
\[
s_c(z)=\frac{15-10}{10}=0.5.
\]

Eso quiere decir que ChatGPT produce 50\% más que el humano en esa tarea.

Como \(a_c(z)>h_c(z)\), el numerador \(a_c(z)-h_c(z)\) es positivo. Como \(h_c(z)>0\), se cumple:
\[
s_c(z)>0.
\]

\subsection{Despeje útil}

La definición:
\[
s_c(z)=\frac{a_c(z)-h_c(z)}{h_c(z)}
\]
puede reordenarse multiplicando ambos lados por \(h_c(z)\):
\[
h_c(z)s_c(z)=a_c(z)-h_c(z).
\]

Por tanto:
\[
a_c(z)-h_c(z)=h_c(z)s_c(z).
\]

Este despeje será clave para expresar la ganancia individual.

\subsection{Promedio del ahorro relativo}

El ahorro promedio condicional a tareas automatizables es:
\[
\rho_c=\frac{1}{\tau_c}\int_{\Zauto}s_c(z)\dd,
\qquad \tau_c>0.
\]

Se divide entre \(\tau_c\) porque \(\tau_c\) es el tamaño del conjunto sobre el que se promedia.

\nota{La regla es la misma que en un promedio discreto: promedio = suma / cantidad. En continuo: promedio = integral / medida del conjunto.}

\section{Paso 4: ganancia individual por usar IA}

Queremos comparar:
\[
\Delta y_{jc}=y_{jc}(1)-y_{jc}(0).
\]

Es decir:

\begin{quote}
ganancia individual = producción con IA - producción sin IA.
\end{quote}

\subsection{Producción con IA}

Partimos de:
\[
y_{jc}(1)=\int_0^1\max\{h_c(z),a_c(z)\}\dd.
\]

Usamos la partición:
\[
[0,1]=([0,1]\setminus\Zauto)\cup\Zauto.
\]

Entonces la integral se parte:
\[
y_{jc}(1)
=\int_{[0,1]\setminus\Zauto}\max\{h_c(z),a_c(z)\}\dd
+\int_{\Zauto}\max\{h_c(z),a_c(z)\}\dd.
\]

En \([0,1]\setminus\Zauto\), ChatGPT no supera al humano:
\[
a_c(z)\leq h_c(z).
\]
Por tanto:
\[
\max\{h_c(z),a_c(z)\}=h_c(z).
\]

En \(\Zauto\), ChatGPT sí supera al humano:
\[
a_c(z)>h_c(z).
\]
Por tanto:
\[
\max\{h_c(z),a_c(z)\}=a_c(z).
\]

Sustituyendo esos dos resultados:
\[
y_{jc}(1)
=\int_{[0,1]\setminus\Zauto}h_c(z)\dd
+\int_{\Zauto}a_c(z)\dd.
\]

\subsection{Producción sin IA}

Sin IA:
\[
y_{jc}(0)=\int_0^1 h_c(z)\dd.
\]

Partimos otra vez el intervalo en las dos mismas regiones:
\[
y_{jc}(0)
=\int_{[0,1]\setminus\Zauto}h_c(z)\dd
+\int_{\Zauto}h_c(z)\dd.
\]

\nota{Con IA, en tareas automatizables aparece \(a_c(z)\). Sin IA, incluso en tareas automatizables aparece \(h_c(z)\), porque no se puede usar ChatGPT.}

\subsection{Resta: lo que se cancela}

Ahora restamos:
\[
\Delta y_{jc}=y_{jc}(1)-y_{jc}(0).
\]

Sustituimos las dos expresiones:
\[
\Delta y_{jc}
=
\left[
\int_{[0,1]\setminus\Zauto}h_c(z)\dd
+\int_{\Zauto}a_c(z)\dd
\right]
-
\left[
\int_{[0,1]\setminus\Zauto}h_c(z)\dd
+\int_{\Zauto}h_c(z)\dd
\right].
\]

Distribuimos el signo menos:
\[
\Delta y_{jc}
=
\int_{[0,1]\setminus\Zauto}h_c(z)\dd
+\int_{\Zauto}a_c(z)\dd
-\int_{[0,1]\setminus\Zauto}h_c(z)\dd
-\int_{\Zauto}h_c(z)\dd.
\]

Los términos sobre tareas no automatizables se cancelan:
\[
\int_{[0,1]\setminus\Zauto}h_c(z)\dd
-\int_{[0,1]\setminus\Zauto}h_c(z)\dd
=0.
\]

Queda:
\[
\Delta y_{jc}
=
\int_{\Zauto}a_c(z)\dd
-\int_{\Zauto}h_c(z)\dd.
\]

Como las dos integrales son sobre el mismo conjunto, podemos unirlas:
\[
\Delta y_{jc}
=
\int_{\Zauto}\left[a_c(z)-h_c(z)\right]\dd.
\]

Usando el despeje del paso anterior:
\[
a_c(z)-h_c(z)=h_c(z)s_c(z),
\]
obtenemos:
\[
\Delta y_{jc}
=
\int_{\Zauto}h_c(z)s_c(z)\dd.
\]

\resultado{Hasta aquí no se impuso ningún supuesto no trivial. La ganancia individual exacta es la suma, sobre tareas automatizables, de productividad humana base multiplicada por ahorro relativo de ChatGPT.}

\section{Paso 4.5: identidad de covarianza}

Ahora queremos convertir:
\[
\Delta y_{jc}=\int_{\Zauto}h_c(z)s_c(z)\dd
\]
en una expresión con promedios.

\subsection{Esperanza condicional dentro de \texorpdfstring{\(\Zauto\)}{Zauto}}

Para cualquier función \(g(z)\), definimos el promedio dentro de \(\Zauto\) como:
\[
\Eauto[g(z)]
=
\frac{1}{\tau_c}\int_{\Zauto}g(z)\dd.
\]

Esto implica:
\[
\int_{\Zauto}g(z)\dd=\tau_c\,\Eauto[g(z)].
\]

Si usamos \(g(z)=h_c(z)s_c(z)\), tenemos:
\[
\int_{\Zauto}h_c(z)s_c(z)\dd
=
\tau_c\,\Eauto[h_c s_c].
\]

Como la integral de la izquierda es \(\Delta y_{jc}\):
\[
\Delta y_{jc}
=
\tau_c\,\Eauto[h_c s_c].
\]

\subsection{Promedios de productividad y ahorro}

El promedio de productividad humana dentro de tareas automatizables es:
\[
\bar h_c
=
\frac{1}{\tau_c}\int_{\Zauto}h_c(z)\dd
=
\Eauto[h_c].
\]

El promedio de ahorro relativo dentro de tareas automatizables es:
\[
\rho_c
=
\frac{1}{\tau_c}\int_{\Zauto}s_c(z)\dd
=
\Eauto[s_c].
\]

\subsection{Derivación de la identidad de covarianza}

La covarianza condicional dentro de \(\Zauto\) se define como:
\[
\Covauto(h_c,s_c)
=
\Eauto\left[
\left(h_c-\Eauto[h_c]\right)
\left(s_c-\Eauto[s_c]\right)
\right].
\]

Expandimos el producto:
\[
\left(h_c-\Eauto[h_c]\right)
\left(s_c-\Eauto[s_c]\right)
=
h_c s_c
-h_c\Eauto[s_c]
-s_c\Eauto[h_c]
+\Eauto[h_c]\Eauto[s_c].
\]

Tomamos esperanza condicional:
\begin{align*}
\Covauto(h_c,s_c)
&=
\Eauto[h_c s_c]
-\Eauto[h_c\Eauto[s_c]]\\
&\quad
-\Eauto[s_c\Eauto[h_c]]
+\Eauto[\Eauto[h_c]\Eauto[s_c]].
\end{align*}

Los términos \(\Eauto[h_c]\) y \(\Eauto[s_c]\) son constantes. Por eso pueden salir de la esperanza:
\[
\Eauto[h_c\Eauto[s_c]]
=
\Eauto[s_c]\Eauto[h_c],
\]
\[
\Eauto[s_c\Eauto[h_c]]
=
\Eauto[h_c]\Eauto[s_c],
\]
\[
\Eauto[\Eauto[h_c]\Eauto[s_c]]
=
\Eauto[h_c]\Eauto[s_c].
\]

Entonces:
\begin{align*}
\Covauto(h_c,s_c)
&=
\Eauto[h_c s_c]
-\Eauto[h_c]\Eauto[s_c]\\
&\quad
-\Eauto[h_c]\Eauto[s_c]
+\Eauto[h_c]\Eauto[s_c]\\
&=
\Eauto[h_c s_c]
-\Eauto[h_c]\Eauto[s_c].
\end{align*}

Reordenamos:
\[
\Eauto[h_c s_c]
=
\Eauto[h_c]\Eauto[s_c]
+\Covauto(h_c,s_c).
\]

Como:
\[
\Eauto[h_c]=\bar h_c,
\qquad
\Eauto[s_c]=\rho_c,
\]
se obtiene:
\[
\Eauto[h_c s_c]
=
\bar h_c\rho_c+\Covauto(h_c,s_c).
\]

\subsection{Sustitución en la ganancia individual}

Antes teníamos:
\[
\Delta y_{jc}
=
\tau_c\,\Eauto[h_c s_c].
\]

Sustituimos la identidad:
\[
\Delta y_{jc}
=
\tau_c\left[
\bar h_c\rho_c+\Covauto(h_c,s_c)
\right].
\]

\resultado{La forma exacta usando covarianza es
\[
\Delta y_{jc}
=
\tau_c\left[
\bar h_c\rho_c+\Covauto(h_c,s_c)
\right].
\]
El factor \(\tau_c\) multiplica tanto al producto de promedios como a la covarianza, porque la covarianza está definida como promedio condicional dentro de \(\Zauto\).}

\subsection{Interpretación de la covarianza}

La covarianza mide si la productividad humana \(h_c(z)\) y el ahorro relativo \(s_c(z)\) se mueven juntos dentro de las tareas automatizables.

\begin{itemize}
  \item Si \(\Covauto(h_c,s_c)>0\), ChatGPT mejora relativamente más tareas donde la productividad humana base era alta.
  \item Si \(\Covauto(h_c,s_c)<0\), ChatGPT mejora relativamente más tareas donde la productividad humana base era baja.
  \item Si \(\Covauto(h_c,s_c)=0\), no hay relación sistemática entre productividad humana y ahorro relativo dentro de tareas automatizables.
\end{itemize}

\supuesto{Para simplificar, el modelo impone \(\Covauto(h_c,s_c)=0\). Entonces:
\[
\Delta y_{jc}=\tau_c\bar h_c\rho_c.
\]}

\section{Paso 4.6: normalización de escala}

El modelo elige la unidad de producción de forma que:
\[
\bar h_c=1.
\]

Esto es una normalización. No significa que empíricamente la productividad humana promedio sea literalmente 1 en unidades naturales. Significa que medimos la producción en unidades donde ese promedio vale 1.

Partimos de:
\[
\Delta y_{jc}=\tau_c\bar h_c\rho_c.
\]

Sustituimos \(\bar h_c=1\):
\[
\Delta y_{jc}=\tau_c(1)\rho_c.
\]

Por tanto:
\[
\Delta y_{jc}=\tau_c\rho_c.
\]

Se define la exposición del lenguaje:
\[
\theta_c=\tau_c\rho_c.
\]

Entonces:
\[
\Delta y_{jc}=\theta_c.
\]

\nota{La exposición del lenguaje combina dos cosas: cuántas tareas del lenguaje son automatizables \((\tau_c)\) y qué tan grande es la mejora promedio en esas tareas \((\rho_c)\).}

\section{Paso 5: decisión individual de entrada al mercado}

Ahora pasamos de productividad a decisión ocupacional.

La utilidad de programar se define como:
\[
\Uprog=\eta+y_{jc}(\mu_j).
\]

Cada componente significa:

\begin{itemize}
  \item \(\eta\): habilidad individual, talento, experiencia o componente personal que aumenta la utilidad/productividad de programar.
  \item \(y_{jc}(\mu_j)\): producción asociada a programar en lenguaje \(c\), con o sin ChatGPT.
\end{itemize}

La utilidad fuera de programar es:
\[
\Uout=r_i.
\]

Aquí \(r_i\) es la utilidad de reserva en el país \(i\). Puede representar salario alternativo, oportunidades fuera de programación, costo de oportunidad u otra actividad.

La persona entra a programar si:
\[
\Uprog\geq\Uout.
\]

Sustituyendo:
\[
\eta+y_{jc}(\mu_j)\geq r_i.
\]

Despejamos \(\eta\). Restamos \(y_{jc}(\mu_j)\) a ambos lados:
\[
\eta\geq r_i-y_{jc}(\mu_j).
\]

El lado derecho es el umbral mínimo de habilidad para entrar:
\[
\eta_c^*(\mu)=r_i-y_{jc}(\mu).
\]

Por tanto:
\[
\text{entra si}\qquad \eta\geq \eta_c^*(\mu).
\]

\subsection{Comparación sin IA y con IA}

Sin IA:
\[
\eta_c^*(0)=r_i-y_{jc}(0).
\]

Con IA:
\[
\eta_c^*(1)=r_i-y_{jc}(1).
\]

Sabemos que:
\[
\Delta y_{jc}=y_{jc}(1)-y_{jc}(0).
\]

Entonces:
\[
y_{jc}(1)=y_{jc}(0)+\Delta y_{jc}.
\]

Sustituimos en el umbral con IA:
\[
\eta_c^*(1)=r_i-\left[y_{jc}(0)+\Delta y_{jc}\right].
\]

Distribuimos el signo negativo:
\[
\eta_c^*(1)=r_i-y_{jc}(0)-\Delta y_{jc}.
\]

Pero:
\[
r_i-y_{jc}(0)=\eta_c^*(0).
\]

Por tanto:
\[
\eta_c^*(1)=\eta_c^*(0)-\Delta y_{jc}.
\]

Usando la versión simplificada:
\[
\Delta y_{jc}=\tau_c\rho_c,
\]
obtenemos:
\[
\eta_c^*(1)=\eta_c^*(0)-\tau_c\rho_c.
\]

Usando la versión exacta con covarianza:
\[
\eta_c^*(1)
=
\eta_c^*(0)
-\tau_c\left[\bar h_c\rho_c+\Covauto(h_c,s_c)\right].
\]

\resultado{El acceso a IA reduce el umbral de entrada. Personas con habilidad entre el nuevo umbral y el antiguo umbral se vuelven programadores marginales que antes no entraban.}

\section{Paso 6: agregación a desarrolladores activos}

Ahora queremos pasar de una persona a una población.

Sea \(N_i\) la población del país \(i\). Sea \(F_i(\eta)\) la función de distribución acumulada de habilidad en el país \(i\):
\[
F_i(x)=\Pr(\eta\leq x).
\]

Esto significa:

\begin{quote}
\(F_i(x)\) es la proporción de personas cuya habilidad es menor o igual que \(x\).
\end{quote}

Entonces:
\[
1-F_i(x)=\Pr(\eta>x).
\]

Como una persona entra si:
\[
\eta\geq \eta_c^*(\mu),
\]
la proporción de personas que entra es aproximadamente:
\[
1-F_i(\eta_c^*(\mu)).
\]

Multiplicando por la población:
\[
\Devs_{ic}(\mu)=N_i\left[1-F_i(\eta_c^*(\mu))\right].
\]

\subsection{Diferencia entre con IA y sin IA}

Con IA:
\[
\Devs_{ic}(1)=N_i\left[1-F_i(\eta_c^*(1))\right].
\]

Sin IA:
\[
\Devs_{ic}(0)=N_i\left[1-F_i(\eta_c^*(0))\right].
\]

Restamos:
\begin{align*}
\Devs_{ic}(1)-\Devs_{ic}(0)
&=
N_i\left[1-F_i(\eta_c^*(1))\right]
-
N_i\left[1-F_i(\eta_c^*(0))\right]\\
&=
N_i\left(
\left[1-F_i(\eta_c^*(1))\right]
-
\left[1-F_i(\eta_c^*(0))\right]
\right)\\
&=
N_i\left[
1-F_i(\eta_c^*(1))-1+F_i(\eta_c^*(0))
\right]\\
&=
N_i\left[
F_i(\eta_c^*(0))-F_i(\eta_c^*(1))
\right].
\end{align*}

Como:
\[
\eta_c^*(1)=\eta_c^*(0)-\tau_c\rho_c,
\]
definimos:
\[
\Delta\eta=\tau_c\rho_c.
\]

Entonces:
\[
\eta_c^*(1)=\eta_c^*(0)-\Delta\eta.
\]

Sustituyendo:
\[
\Devs_{ic}(1)-\Devs_{ic}(0)
=
N_i\left[
F_i(\eta_c^*(0))
-F_i(\eta_c^*(0)-\Delta\eta)
\right].
\]

\subsection{Aproximación de Taylor}

Sea:
\[
x=\eta_c^*(0).
\]

Queremos aproximar:
\[
F_i(x)-F_i(x-\Delta\eta).
\]

La expansión de Taylor de primer orden de \(F_i\) alrededor de \(x\) es:
\[
F_i(x+\Delta)\approx F_i(x)+F_i'(x)\Delta.
\]

Aquí queremos \(x-\Delta\eta\), así que usamos \(\Delta=-\Delta\eta\):
\[
F_i(x-\Delta\eta)
\approx
F_i(x)+F_i'(x)(-\Delta\eta).
\]

Por tanto:
\[
F_i(x-\Delta\eta)
\approx
F_i(x)-F_i'(x)\Delta\eta.
\]

La densidad \(f_i\) es la derivada de la CDF:
\[
f_i(x)=F_i'(x).
\]

Entonces:
\[
F_i(x-\Delta\eta)
\approx
F_i(x)-f_i(x)\Delta\eta.
\]

Ahora:
\begin{align*}
F_i(x)-F_i(x-\Delta\eta)
&\approx
F_i(x)-\left[F_i(x)-f_i(x)\Delta\eta\right]\\
&=
F_i(x)-F_i(x)+f_i(x)\Delta\eta\\
&=
f_i(x)\Delta\eta.
\end{align*}

Volvemos a \(x=\eta_c^*(0)\):
\[
F_i(\eta_c^*(0))
-F_i(\eta_c^*(0)-\Delta\eta)
\approx
f_i(\eta_c^*(0))\Delta\eta.
\]

Como \(\Delta\eta=\tau_c\rho_c\):
\[
F_i(\eta_c^*(0))
-F_i(\eta_c^*(0)-\Delta\eta)
\approx
f_i(\eta_c^*(0))\tau_c\rho_c.
\]

Sustituyendo en la diferencia de desarrolladores:
\[
\Devs_{ic}(1)-\Devs_{ic}(0)
\approx
N_i f_i(\eta_c^*(0))\tau_c\rho_c.
\]

\supuesto{Esta aproximación de Taylor es buena si \(\Delta\eta=\tau_c\rho_c\) es pequeño respecto a la escala en la que cambia la densidad \(f_i\). Si el cambio de umbral es grande, los términos de segundo orden podrían importar.}

\nota{El término \(f_i(\eta_c^*(0))\) mide cuánta gente hay cerca del umbral inicial. Si hay mucha gente justo debajo del umbral, una pequeña reducción del umbral genera muchos nuevos desarrolladores.}

\section{Paso 7: resultado per cápita, absorción y separabilidad}

El resultado anterior está en número total de desarrolladores:
\[
\Devs_{ic}(1)-\Devs_{ic}(0)
\approx
N_i f_i(\eta_c^*(0))\tau_c\rho_c.
\]

Pero el outcome del estudio se expresa como desarrolladores por cada 100,000 habitantes. Por eso se define:
\[
\Delta Y_{ic}
=
\frac{\Devs_{ic}(1)-\Devs_{ic}(0)}{N_i}\cdot 10^5.
\]

Como \(10^5=100{,}000\), esta fórmula dice:

\begin{quote}
Tomamos el cambio total de desarrolladores, lo dividimos entre la población y lo multiplicamos por 100,000.
\end{quote}

\subsection{Sustitución paso a paso}

Partimos de:
\[
\Delta Y_{ic}
=
\frac{\Devs_{ic}(1)-\Devs_{ic}(0)}{N_i}\cdot 10^5.
\]

Sustituimos:
\[
\Devs_{ic}(1)-\Devs_{ic}(0)
\approx
N_i f_i(\eta_c^*(0))\tau_c\rho_c.
\]

Entonces:
\[
\Delta Y_{ic}
\approx
\frac{
N_i f_i(\eta_c^*(0))\tau_c\rho_c
}{N_i}
\cdot 10^5.
\]

Cancelamos \(N_i\):
\[
\Delta Y_{ic}
\approx
f_i(\eta_c^*(0))\tau_c\rho_c\cdot 10^5.
\]

Reordenamos:
\[
\Delta Y_{ic}
\approx
10^5 f_i(\eta_c^*(0))\tau_c\rho_c.
\]

\subsection{Factorización en país y lenguaje}

La expresión:
\[
10^5 f_i(\eta_c^*(0))\tau_c\rho_c
\]
tiene dos partes naturales.

Primera parte:
\[
10^5 f_i(\eta_c^*(0)).
\]

Esta parte depende del país porque \(f_i\) es la distribución de habilidad del país \(i\). También mide cuánta gente hay cerca del umbral de entrada. Se interpreta como capacidad de absorción.

Segunda parte:
\[
\tau_c\rho_c.
\]

Esta parte depende del lenguaje \(c\). Mide cuántas tareas del lenguaje son automatizables y cuánto mejora ChatGPT en promedio.

Por eso se define:
\[
\theta_c=\tau_c\rho_c.
\]

Y, bajo un supuesto de separabilidad, se define:
\[
\phi_i=10^5 f_i(\eta_c^*(0)).
\]

Entonces:
\[
\Delta Y_{ic}\approx \phi_i\theta_c.
\]

\resultado{El aumento per cápita se escribe como producto entre exposición del lenguaje y absorción del país:
\[
\Delta Y_{ic}\approx \phi_i\theta_c.
\]}

\subsection{Por qué la separabilidad es un supuesto}

El umbral inicial era:
\[
\eta_c^*(0)=r_i-y_{jc}(0).
\]

Aquí:

\begin{itemize}
  \item \(r_i\) depende del país \(i\).
  \item \(y_{jc}(0)\) depende del lenguaje \(c\), porque la productividad base puede ser distinta entre lenguajes.
\end{itemize}

Por tanto, en rigor:
\[
f_i(\eta_c^*(0))
\]
también puede depender de \(c\). Si el umbral cambia por lenguaje, la densidad evaluada en ese umbral también cambia por lenguaje.

La forma más general sería:
\[
\phi_{ic}=10^5 f_i(\eta_c^*(0)).
\]

Entonces:
\[
\Delta Y_{ic}\approx \phi_{ic}\theta_c.
\]

Para escribir:
\[
\Delta Y_{ic}\approx \phi_i\theta_c,
\]
se necesita asumir que:
\[
\phi_{ic}\approx \phi_i.
\]

Esto significa:

\begin{quote}
La capacidad de absorción depende principalmente del país y no cambia mucho entre lenguajes.
\end{quote}

Una manera de justificarlo es asumir que las productividades base \(y_{jc}(0)\) son similares entre lenguajes, de modo que los umbrales \(\eta_c^*(0)\) no se mueven demasiado. Otra manera es asumir que la densidad \(f_i\) es relativamente plana alrededor de esos umbrales, de modo que evaluarla en umbrales ligeramente distintos da valores parecidos.

\supuesto{La separabilidad país--lenguaje permite reemplazar \(\phi_{ic}\) por \(\phi_i\). Sin este supuesto, el efecto sigue pudiendo interpretarse por lenguaje, pero la absorción promedio podría depender también del lenguaje.}

\section{Paso 8: especificación de panel}

Finalmente se lleva el resultado teórico a una ecuación empírica. La idea teórica era:
\[
\Delta Y_{ic}\approx \phi_i\theta_c.
\]

Pero ese cambio solo ocurre cuando ChatGPT está disponible. Sea:
\[
A_{it}\in\{0,1\}
\]
un indicador de disponibilidad de ChatGPT en el país \(i\) y periodo \(t\).

Entonces el componente de tratamiento es:
\[
\theta_c\phi_i A_{it}.
\]

Si \(A_{it}=0\), no hay tratamiento:
\[
\theta_c\phi_i A_{it}=0.
\]

Si \(A_{it}=1\), el efecto esperado es:
\[
\theta_c\phi_i A_{it}=\theta_c\phi_i.
\]

La línea base sin tratamiento puede descomponerse como:
\[
Y_{ic}(0)=\alpha_i+\beta_t+\delta_c+\eta_{ict}.
\]

Donde:

\begin{itemize}
  \item \(\alpha_i\): diferencias permanentes entre países.
  \item \(\beta_t\): shocks comunes de tiempo.
  \item \(\delta_c\): diferencias permanentes entre lenguajes.
  \item \(\eta_{ict}\): factores no observados restantes.
\end{itemize}

Sumando el componente de tratamiento:
\[
Y_{ict}
=
\alpha_i+\beta_t+\delta_c+\theta_c\phi_i A_{it}+\eta_{ict}.
\]

\subsection{Condición para interpretación causal}

Para interpretar causalmente el coeficiente de tratamiento, se requiere:
\[
\mathbb{E}\left[\eta_{ict}\mid \alpha_i,\beta_t,\delta_c,A_{it}\right]=0.
\]

Esto significa:

\begin{quote}
Después de controlar por efectos fijos de país, tiempo y lenguaje, la disponibilidad de ChatGPT no debe estar correlacionada con shocks no observados que también afecten \(Y_{ict}\).
\end{quote}

Si esta condición falla, parte del cambio atribuido a ChatGPT podría venir de otros factores.

\section{Resumen de la derivación completa}

\begin{center}
\begin{tabular}{p{0.08\linewidth}p{0.52\linewidth}p{0.29\linewidth}}
\toprule
Paso & Ecuación clave & Tipo \\
\midrule
1 &
\(\displaystyle y_{jc}(\mu_j)=\int_0^1\max\{h_c(z),\mu_j a_c(z)\}\dd\) &
Definición \\
2 &
\(\displaystyle \Zauto=\{z:a_c(z)>h_c(z)\}\), \(\displaystyle \tau_c=\int_{\Zauto}1\dd\) &
Definición exacta \\
3 &
\(\displaystyle s_c(z)=\frac{a_c(z)-h_c(z)}{h_c(z)}\), \(\displaystyle \rho_c=\frac{1}{\tau_c}\int_{\Zauto}s_c(z)\dd\) &
Definición exacta \\
4 &
\(\displaystyle \Delta y_{jc}=\int_{\Zauto}h_c(z)s_c(z)\dd\) &
Identidad exacta \\
4.5 &
\(\displaystyle \Delta y_{jc}=\tau_c[\bar h_c\rho_c+\Covauto(h_c,s_c)]\) &
Identidad de covarianza \\
4.6 &
\(\displaystyle \Delta y_{jc}=\tau_c\rho_c=\theta_c\) &
Covarianza cero + normalización \\
5 &
\(\displaystyle \eta_c^*(1)=\eta_c^*(0)-\tau_c\rho_c\) &
Decisión de entrada \\
6 &
\(\displaystyle \Devs_{ic}(1)-\Devs_{ic}(0)\approx N_i f_i(\eta_c^*(0))\tau_c\rho_c\) &
Taylor de primer orden \\
7 &
\(\displaystyle \Delta Y_{ic}\approx \phi_i\theta_c\) &
Per cápita + separabilidad \\
8 &
\(\displaystyle Y_{ict}=\alpha_i+\beta_t+\delta_c+\theta_c\phi_i A_{it}+\eta_{ict}\) &
Ecuación de panel \\
\bottomrule
\end{tabular}
\end{center}

\section{Mini-ejemplo numérico}

Supongamos que en un lenguaje:
\[
\tau_c=0.30,\qquad \rho_c=0.50.
\]

Esto significa:

\begin{itemize}
  \item ChatGPT supera al humano en 30\% de las tareas.
  \item En esas tareas, la mejora relativa promedio es 50\%.
\end{itemize}

Entonces:
\[
\theta_c=\tau_c\rho_c=0.30\times 0.50=0.15.
\]

La ganancia individual simplificada es:
\[
\Delta y_{jc}=0.15.
\]

Si el umbral inicial era:
\[
\eta_c^*(0)=2.00,
\]
el nuevo umbral es:
\[
\eta_c^*(1)=2.00-0.15=1.85.
\]

Las personas con habilidad:
\[
1.85\leq \eta<2.00
\]
no entraban antes, pero entran después.

Si el país tiene:
\[
N_i=10{,}000{,}000,
\]
y la densidad alrededor del umbral inicial es:
\[
f_i(\eta_c^*(0))=0.08,
\]
entonces:
\[
\Devs_{ic}(1)-\Devs_{ic}(0)
\approx
10{,}000{,}000\times 0.08\times 0.15
=120{,}000.
\]

Por cada 100,000 habitantes:
\[
\Delta Y_{ic}
=
\frac{120{,}000}{10{,}000{,}000}\times 100{,}000
=1{,}200.
\]

\nota{El ejemplo es puramente ilustrativo. Sirve para ver cómo se mueven las piezas, no para afirmar magnitudes reales.}

\section{Lista de supuestos importantes}

\begin{enumerate}
  \item \textbf{Elección del mejor método por tarea.} El modelo asume que el desarrollador puede usar el método más productivo en cada tarea.
  \item \textbf{Covarianza cero.} Para obtener \(\Delta y_{jc}=\tau_c\rho_c\bar h_c\), se impone \(\Covauto(h_c,s_c)=0\).
  \item \textbf{Normalización de escala.} Se elige \(\bar h_c=1\), lo que permite \(\Delta y_{jc}=\tau_c\rho_c\).
  \item \textbf{Taylor de primer orden.} Se aproxima el cambio en desarrolladores usando la densidad en el umbral inicial.
  \item \textbf{Separabilidad país--lenguaje.} Se reemplaza \(\phi_{ic}\) por \(\phi_i\), asumiendo que la absorción depende principalmente del país.
  \item \textbf{Exogeneidad condicional.} Para interpretación causal en el panel, se requiere que la disponibilidad de ChatGPT no esté correlacionada con shocks no observados después de controlar por efectos fijos.
\end{enumerate}

\section{Lectura intuitiva final}

El modelo dice:

\[
\underbrace{\text{más tareas automatizables}}_{\tau_c}
\times
\underbrace{\text{mayor ahorro promedio}}_{\rho_c}
\Longrightarrow
\underbrace{\text{mayor ganancia individual}}_{\Delta y_{jc}}.
\]

Luego:
\[
\underbrace{\text{mayor ganancia individual}}_{\Delta y_{jc}}
\Longrightarrow
\underbrace{\text{menor umbral de entrada}}_{\eta_c^*(1)<\eta_c^*(0)}.
\]

Luego:
\[
\underbrace{\text{menor umbral}}_{\text{baja en }\eta_c^*}
\times
\underbrace{\text{gente cerca del umbral}}_{f_i(\eta_c^*(0))}
\Longrightarrow
\underbrace{\text{más desarrolladores activos}}_{\Devs_{ic}(1)-\Devs_{ic}(0)}.
\]

Finalmente:
\[
\Delta Y_{ic}\approx \phi_i\theta_c.
\]

En una frase:

\begin{quote}
ChatGPT tiene mayor impacto donde los lenguajes están más expuestos a automatización y donde los países tienen más personas cerca del umbral de entrada a la programación.
\end{quote}

\end{document}
